% Independent manuscript fragment. Requires amsmath, amssymb, amsthm.
\subsection{Candidate-local simulation in the native experiment}
\label{sec:candidate-local-accounting}

Which record can the reduction leave opaque? The answer depends on
which two reservations supplied the accepted send. We therefore localize
the score in the native experiment first: each eligible send with a unique
usable pair is assigned to that pair, ineligible runs contribute zero,
and eligible attribution failures remain in the original game. We can then
compare each candidate using the records, protected report rows
and stopping rule determined by its own reservations.

Let $\mathsf G_{\rm nat}$ be the full native experiment and
$J=[G]\times[C]$ its fixed reservation-slot budget.  Write $E$ for original
eligibility, $A$ for usable unique attribution, and $I\in J$ for the attributed
pair on $A$.  These variables describe the actual accepted logical challenge
send; selection does not transfer its test bit to another send.  For a bounded
application score $S\in[-1,1]$, set
\[
 Z=\mathbf1_E S,\qquad
 Z_j=\mathbf1_{E\land A\land(I=j)}S.
\]

For each $j$, let $U_j$ and $B_j$ be efficient, cumulative prefix events
computed from the available classical provenance ledger.  They are
noninterfering monitors: before stopping they change no native response or
protocol transition and do not inspect an unknown tape or a quantum query.
A \emph{free stop} $U_j$ permanently falsifies $E\land A\land(I=j)$ for every
continuation.  A \emph{charged stop} $B_j$ may instead be a permitted
protected-report-key projection.  Define $\mathsf G_{\rm nat}^{[j]}$ to run
the native code until the first such event, then enter a stopped state with
candidate score zero before returning the monitored bytes.  An event for one
candidate does not stop the original game or another candidate's reduction.

\begin{lemma}[Origin-anchored adaptive stopping]
\label{lem:adaptive-origin-stop}
Fix candidate $j$ and one of its original reservation slots $o$ which has
occurred.  Suppose a successful capture supplies an intact version $v$ with
\[
 \mathsf{subject}(v)=\mathsf{session}(o),\qquad
 \mathsf{res}(o)\le\mathsf{save}(v).
\]
Assume usable attribution to $j$ assigns its member origins to the selected
session, and eligibility places the cut before both reservations.  Then this
capture is a free $U_j$ event, even if the other reservation or challenge send
has not occurred and the final cut has not been selected.
\end{lemma}
\begin{proof}
In any continuation with $E\land A\land(I=j)$, the selected session is
$s=\mathsf{session}(o)$ and the selected cut satisfies
$c<\mathsf{res}(o)\le\mathsf{save}(v)$.  Acquisition has already inserted
$v$ into the cumulative archive. Thus the execution has already acquired a
version of $s$ saved after $c$, contradicting the original acquisition conjunct.
Copies, resealing, deletion, and recovery preserve this contradiction because
the subject, original save event, and acquisition are immutable.  If the other
reservation or challenge send never occurs, the attributed eligible event is
already false.  Decryption and a final cut value are unnecessary for this
certificate.
\end{proof}

Actual retention of $o$ in $v$ supplies the time inequality by reservation/save
causality; an independently recorded later save of the same session also
suffices.  More generally, let $L_s$ be the largest original save time of an
acquired intact version of session $s$, with $L_s=-\infty$ for an empty set.
For an existing pair $(g,a)$, the original temporal eligibility conditions are
exactly
\[
       L_s\le c<\min\{\mathsf{res}(g),\mathsf{res}(a)\}.
\]
Once the interval is empty, the selector can no longer choose an
eligible cut for this pair. Resealing the acquired version preserves
its first-save time and hence $L_s$. A batch-only key fragment or a
public commitment can depend on the target session without containing
an intact version. Neither raises $L_s$ on that account. A cut inside
the interval must also meet the remaining eligibility conditions.

\paragraph{One native residual.}
In the full native run let $B_I$ denote the charged event for its final
attributed pair, used only on $A$, and define
\[
 \eta_{\rm attr}^{\rm nat}=\Pr_{\mathsf G_{\rm nat}}[E\land\neg A],
 \qquad
 \eta_{\rm proj}^{\rm nat}=\Pr_{\mathsf G_{\rm nat}}[E\land A\land B_I].
\]
The two residuals measure disjoint events in one game. In particular,
\[
 \eta_{\rm proj}^{\rm nat}
 =\sum_{j\in J}\Pr_{\mathsf G_{\rm nat}}
       [E\land A\land(I=j)\land B_j].
\]
The projection residual counts the charged event for the selected pair. For
the report construction, $B_j$ is the actual protected-key projection event
for the closure induced by that candidate's source family. A projection
affecting only another candidate therefore contributes nothing to this
residual. The selected projection probability remains explicit in the
native-game bound, and the experiment's original session-acquisition rule
is preserved.

\begin{lemma}[Native score localization with local stops]
\label{lem:native-local-score}
Let $\Psi_j$ be the same terminal code for
$\mathbf1_{E\land A\land(I=j)}S$ in each experiment, with value zero in its
stopped state.  Then
\begin{align}
 \mathbb E_{\mathsf G_{\rm nat}^{[j]}}\Psi_j
   &=\mathbb E_{\mathsf G_{\rm nat}}
           [Z_j\mathbf1_{\neg B_j}], \label{eq:local-stop-identity}\\
 \mathbb E_{\mathsf G_{\rm nat}}Z
   &=\mathbb E_{\mathsf G_{\rm nat}}[Z\mathbf1_{\neg A}]
     +\mathbb E_{\mathsf G_{\rm nat}}[Z\mathbf1_{A\land B_I}]
     +\sum_{j\in J}
       \mathbb E_{\mathsf G_{\rm nat}^{[j]}}\Psi_j.
       \label{eq:native-local-decomposition}
\end{align}
\end{lemma}
\begin{proof}
Until its first local stop, the stopped machine has the native prefix
distribution.  On a run with $B_j$, its candidate score is zero.  On a run
with $U_j$, every native continuation already has $Z_j=0$, whether or not
$B_j$ later occurs.  On all other runs the executions and candidate scores
agree.  This proves~\eqref{eq:local-stop-identity}.  Partition the one native
score by attribution and its unique index; pointwise,
$\sum_j Z_j\mathbf1_{B_j}=Z\mathbf1_{A\land B_I}$.  Substitution proves
\eqref{eq:native-local-decomposition}. This identity uses the final index
only to attribute the completed run. For a quantum adversary the
argument compares identical classical-prefix instruments and their stopped
continuations, retaining the native quantum-state transformations and
oracle-query interface up to the local stop.
\end{proof}

\paragraph{Local cryptographic comparisons.}
For each $j$, let $\mathsf V^{[j]}$ be its specified real-component
continuation after the local source/record/report replacements, and let
$\widetilde{\mathsf V}^{[j]}$ be its corresponding random-component game.
All terminal events are recomputed in their own game by the same score code;
no event from a native transcript is transferred to an altered transcript.
Suppose the actual stopped-interface correspondence gives
\[
 \left|\mathbb E_{\mathsf G_{\rm nat}^{[j]}}\Psi_j
       -\mathbb E_{\mathsf V^{[j]}}\Psi_j\right|\le\rho_j.
\]
This is a local score bound, including its native source correspondence.
A binary record or report bound contributes twice its value to this score
comparison.  The stopping certificate does not itself establish the
correspondence for permitted partial projections.

For the remaining component step, specialize to $S\in\{-1,+1\}$, so that
$\Psi_j\in\{-1,0,+1\}$.  Suppose its actual fixed-pair adapter satisfies,
for $\sigma\in\{-1,+1\}$,
\[
 \left|\Pr_{\mathsf V^{[j]}}[\Psi_j=\sigma]
       -\Pr_{\widetilde{\mathsf V}^{[j]}}[\Psi_j=\sigma]\right|
 \le a_j^\sigma+\delta_j.
\]
Here $a_j^\sigma$ is evaluated for the complete derived distinguisher,
including its shared-oracle work; $\delta_j$ is the selected pair's correctness term. Write
$b_j=\mathbb E_{\widetilde{\mathsf V}^{[j]}}\Psi_j$.

\begin{proposition}[Candidate-local native accounting]
\label{prop:native-local-accounting}
Under the stated local comparisons, the native message-guessing advantage
satisfies
\[
 \operatorname{Adv}_{\rm msg}
 \le\min\!\left\{\frac12,\frac12\!\left(
 \eta_{\rm attr}^{\rm nat}+\eta_{\rm proj}^{\rm nat}
 +\sum_{j\in J}
    \bigl(\rho_j+a_j^++a_j^-+2\delta_j+|b_j|\bigr)
 \right)\right\}.
\]
\end{proposition}
\begin{proof}
The first two terms of~\eqref{eq:native-local-decomposition} have magnitudes
at most the two native residuals.  The local correspondence costs $\rho_j$.
Subtracting the positive- and negative-event comparisons bounds the local
component score change by $a_j^++a_j^-+2\delta_j$.  Sum the candidate terms,
apply the triangle inequality, and use $|Z|\le1$.  Completing invalid runs
by an independent fair guess gives
$\operatorname{Adv}_{\rm msg}=\tfrac12|\mathbb E_{\mathsf G_{\rm nat}}Z|$.
\end{proof}

The ideal continuations may have different adaptive selectors, so each
$|b_j|$ belongs to its own continuation and remains a separate term.
If a common intermediate game and a global
comparison have separately been established, that shared error can instead be
charged once before the common-game accounting. The proposition above uses
the local comparisons available for each candidate.
